\unitopen{34}
\begin{candos}
\item zreferować w standupie, co zrobiłeś, co robisz i co Cię blokuje;
\item opisać błąd, jego przyczynę i usunięcie w Perfekcie i Passiv;
\item poprosić o review, zgłosić problem i ustalić termin;
\item przeczytać i napisać krótki opis pull requestu po niemiecku.
\end{candos}

\section{Słownictwo}
\vocabtable{34}

\section{Standup po niemiecku}
Trzy pytania: \de{Was habe ich gestern gemacht?} \de{Was mache ich heute?} \de{Was blockiert mich?} Wzór:
\de{Gestern habe ich den Bug im Login gefunden und behoben. Heute arbeite ich an der API für die Rechnungen. Ich bin blockiert,
weil mir der Zugriff auf die Datenbank fehlt.} Czasowniki angielskie dostają niemieckie końcówki: \de{mergen} (\de{gemergt}),
\de{deployen} (\de{deployt}), \de{committen} (\de{committet}), \de{reviewen} (\de{reviewt}). Dla pisania formalnego użyj \de{zusammenführen}
zamiast \de{mergen}, ale w zespole \de{mergen} jest normą.

\section{Opis błędu i incydentu}
\begin{center}\small
\begin{tabular}{@{}p{0.2\linewidth}p{0.72\linewidth}@{}}
\toprule
\textsf{\textbf{Element}} & \textsf{\textbf{Zdanie}}\\
\midrule
Objaw & \de{Die Anwendung stürzt beim Login ab.}\\
Zasięg & \de{Das Problem betrifft alle Benutzer in Europa.}\\
Przyczyna & \de{Die Ursache ist ein abgelaufenes Zertifikat.}\\
Działanie & \de{Wir haben das Zertifikat erneuert und den Dienst neu gestartet.}\\
Zapobieganie & \de{Damit das nicht wieder passiert, richten wir einen Alarm ein.}\\
\bottomrule
\end{tabular}
\end{center}

\begin{falle}
Polska skłonność do łączenia terminów angielskich z polskimi końcówkami („zmergowałem”) ma niemiecki odpowiednik, ale z regułami:
\de{gemergt}, \de{deployt}; \emph{nie} \de{gemerged}. Drugi błąd: używanie \de{Bug} bez rodzajnika; rodzajnik jest męski: \de{der Bug}.
Trzeci: tłumaczenie „zgłosić” jako \de{zeigen}: \de{Ich habe den Bug gemeldet} (zgłosiłem), nie \de{gezeigt}.
\end{falle}

\begin{szwajcaria}
W szwajcarskich zespołach IT rozmawia się często po angielsku, ale język oficjalny to niemiecki, zwłaszcza w komunikacji z
klientem i dokumentacji. Termin \de{Sitzung} bywa używany zamiast \de{Meeting}. W Zurychu powszechne są firmy międzynarodowe,
w których język zespołu to angielski; niemiecki pomoże Ci jednak w integracji i w kontakcie z biznesem.
\end{szwajcaria}

\begin{praca}
Opis PR: \de{Dieser Pull Request führt einen Cache für die Abfrage der Kundendaten ein. Die Änderung reduziert die Antwortzeit
um etwa 40 Prozent. Bitte prüft besonders die Invalidierung des Caches.} Prośba: \de{Könnte jemand bis Mittwoch reviewen?}
\end{praca}

\section{Dialogi}
\textbf{Dialog 1 --- Daily}
\begin{dialog}
\turn{A}{Guten Morgen zusammen. Wer beginnt?}{Dzień dobry wszystkim. Kto zaczyna?}
\turn{B}{Ich. Gestern habe ich den Bug im Login behoben und einen Pull Request erstellt. Heute reviewe ich den Code von Anna.}{Ja. Wczoraj usunąłem błąd w logowaniu i utworzyłem pull request. Dziś robię review kodu Anny.}
\turn{C}{Ich arbeite heute an der Schnittstelle. Ich bin blockiert, weil ich keinen Zugriff auf die Testumgebung habe.}{Dziś pracuję nad interfejsem. Jestem zablokowany, bo nie mam dostępu do środowiska testowego.}
\turn{A}{Darum kümmere ich mich. Ich schicke dir die Berechtigung bis zehn Uhr.}{Zajmę się tym. Wyślę ci uprawnienie do dziesiątej.}
\end{dialog}

\section{Ćwiczenia}
\begin{cwiczenia}
\exitem{Powiedz, że wczoraj naprawiłeś błąd (Perfekt).}{\de{Gestern habe ich den Bug behoben.}}
\exitem{Powiedz, że dziś pracujesz nad API.}{\de{Heute arbeite ich an der API.}}
\exitem{Powiedz, że jesteś zablokowany z powodu braku dostępu.}{\de{Ich bin blockiert, weil mir der Zugriff fehlt.}}
\exitem{Partizip II: \de{deployen}.}{\de{deployt}}
\exitem{Partizip II: \de{mergen}.}{\de{gemergt}}
\exitem{Passiv Präteritum: \de{Man erneuerte das Zertifikat.}}{\de{Das Zertifikat wurde erneuert.}}
\exitem{Podaj przyczynę awarii (zdanie z \de{Die Ursache ist \dots}).}{np.\ \de{Die Ursache ist ein abgelaufenes Zertifikat.}}
\exitem{Poproś o review do środy.}{\de{Könnte jemand bis Mittwoch reviewen?}}
\exitem{Popraw: \de{Ich habe den Bug gezeigt.} (chodzi o zgłoszenie)}{\de{Ich habe den Bug gemeldet.}}
\exitem{Przetłumacz: „Potrzebuję pomocy przy konfiguracji.”}{\de{Ich brauche Hilfe bei der Konfiguration.}}
\exitem{Przetłumacz: „To zostało załatwione.”}{\de{Das ist erledigt.}}
\end{cwiczenia}

\section{Trening słownictwa (generowany)}
\drill{34}

\weekend{Napisz niemiecki opis jednego prawdziwego PR lub incydentu (objaw, zasięg, przyczyna, działanie, zapobieganie)
w dziesięciu zdaniach i wygłoś trzyminutowy standup po niemiecku. Nagraj i sprawdź czasy.}
